%%%%%%%%%%%%%%%%%%%%%%%%%%%%%%%%%
% Conclusion
%%%%%%%%%%%%%%%%%%%%%%%%%%%%%%%%%

\section{Conclusion}
\label{sec:conclusion}

Under cross-set independence, complement, union, intersection, and set difference of Bernoulli sets yield (possibly higher-order) Bernoulli sets whose per-region error probabilities are closed-form functions of the operand rates, and whose marginal rates are cardinality-weighted averages of them.
These are the standard independent-OR and independent-AND rules~\cite{broderMitzenmacher2004,parker1975,dalvi2007}, collected here in one notation.
At the level of realized sets, union and intersection form commutative monoids with the empty set and universal set as identities.
Relational predicates---subset and equality---admit exact probability formulas that factorize element-wise.
An interval arithmetic extension propagates uncertain rate parameters through set operations, yielding guaranteed bounds when parameters are only approximately known.

\paragraph{Practical implications.}
For a system built on independently constructed Bloom filters, the formulas of \cref{sec:set_theory} give the error rates of any set-theoretic expression, provided they are applied region by region: the per-element probability of each intermediate result must be tracked by the membership pattern of the element in all operands (see the remark on nested expressions in \cref{sec:set_ops_summary}).
Reusing the marginal (FPR, FNR) pair of an intermediate result in a further operation is not valid in general, and confidence intervals for marginal rates must account for the Poisson-binomial error count.

\paragraph{Open problems.}
The partition weights $w_1, w_2, w_3$ in the set-operation formulas require knowledge of the objective set cardinalities, which are typically unknown (\cref{rem:unknown_partitions}).
Developing principled estimators for these cardinalities from the approximate sets themselves, with quantified bias and variance, is a natural next step.

\paragraph{Companion papers.}
The Bernoulli set model, including the axioms, error-count distributions, asymptotic limits, confidence intervals, and higher-order composition theorem, is developed in~\cite{bernoulliSets}.
The probability distributions of binary classification measures (positive predictive value, negative predictive value, accuracy, Youden's $J$) are derived in~\cite{bernoulliMeasures}.
The joint entropy of error counts and the information-theoretic space complexity lower bound are developed in~\cite{bernoulliEntropy}.
